\documentclass[10pt,a4paper]{amsart}
\usepackage[margin=19mm]{geometry}
\usepackage{amsmath,amssymb,mathtools}
\usepackage[scr=rsfs,cal=euler]{mathalfa}
\usepackage{xfrac}
\usepackage[unicode,hidelinks,bookmarks=false]{hyperref}
\newcommand{\ie}{i.e.\ }
\newcommand{\cf}{cf.\ }
\newcommand{\resp}{respectively\ }
\newcommand{\Subsection}{\subsection}
\providecommand{\nobreakdash}{\nobreak\hbox{-}}
\setlength{\emergencystretch}{2em}
\multlinegap0pt
\usepackage{amssymb}
\usepackage{algorithm}\usepackage{algpseudocode}
\usepackage{mathtools}
\usepackage{xcolor}
\newtheorem{definition}{Definition}
\newtheorem{proposition}{Proposition}
\newcommand{\supp}{\mathrm{supp}}
\title{Support Recovery in One-bit Compressed Sensing with Near-Optimal Measurements and Sublinear Time}
\author{Xiaxin Li, Arya Mazumdar}
\date{Source: arXiv:2511.10777v4 (CC BY 4.0)}
\begin{document}
\maketitle

\noindent\emph{Source and licence.} Support Recovery in One-bit Compressed Sensing with Near-Optimal Measurements and Sublinear Time. Version arXiv:2511.10777v4 (CC BY 4.0), distributed by arXiv under the Creative Commons Attribution 4.0 International licence, http://creativecommons.org/licenses/by/4.0/, as recorded in the arXiv licence metadata for this exact version and shown on the abstract page https://arxiv.org/abs/2511.10777v4. Full licence notice: \texttt{licenses/arxiv-2511.10777-CC-BY-4.0.txt}.\par\smallskip

\noindent\textit{Editorial scope.} Counted unit: the article's proposition prop2.1.1 (source0.tex lines 436--438). The article's complete original proof of it is delivered with it (source0.tex lines 439--469, one \texttt{proof} environment of 31 lines). No mathematics is written, repaired or re-derived: the statement and the proof are the article's own lines, in the article's own notation, and no retained line is substituted or re-flowed.  Retained uncounted as the source-local context the proof needs: lines 361--377.  Nothing was omitted from any retained span, so the package carries no omission record.  No retained line contains a citation, so the wrapper carries no bibliography and no bibliography entry.  No retained line contains a figure environment, an image-inclusion command or drawing code, so no image is delivered and the package declares no asset dependency.  Editorial formatting only, in addition: an AMS \texttt{amsart} preamble that restates the article's own declarations and macros used by the retained lines, namely the environments \texttt{definition}, \texttt{proposition} on separate counters, together with the standard packages those lines need.  The retained environments print in this document as Definition 1, Proposition 1 in order of appearance, whereas the article prints the same items with the numbering of its own full document.  The complete native source of the article as distributed in the arXiv e-print for this exact version is bundled unchanged as \texttt{sources/189064/source0.tex}.
\begin{definition}\label{def:kld}
A binary matrix $M \in \{0,1\}^{m \times n}$ is called $(k,l)$-\emph{distinguishable} if for any vector $x \in \{0,1\}^n$ with $|\supp(x)| = k$, we have
% Original version: 
\[
\left| \left\{ i \in \supp(x) : \exists j \in [m] \text{ s.t. } \supp(M^j) \cap \supp(x) = \{i\} \right\} \right| > k - l.
\]

% % IEEE version: 
% \begin{equation*}
% \begin{aligned}
% \Bigl| \Bigl\{ i \in \supp(x) :\; &\exists j \in [m] \text{ s.t. } \\ 
% &\supp(M^j) \cap \supp(x) = \{i\} \Bigr\} \Bigr|
% > k - l.
% \end{aligned}
% \end{equation*}

\end{definition}

\begin{proposition}\label{prop2.1.1}
    A $(n,m,d,k,l, 1)$-list UF matrix is $(k,l)$-distinguishable.
\end{proposition}

\begin{proof}
    Let $M$ be an $(n,m,d,k,l,1)$-list union-free (UF) matrix. Consider any set of $k$ columns of $M$, denoted by $M_{j_1}, \dots, M_{j_k}$. Let $V:=\{j_1,...,j_k\}$, and
\(
M' := M_{V} = [M_{j_1}, \dots, M_{j_k}]
\)
be the corresponding submatrix. 
By Definition~\ref{def:kld}, to show that $M$ is $(k,l)$-distinguishable, it suffices to exhibit a submatrix of $M'$ consisting of at least $k - l + 1$ rows, each equal to a distinct standard basis vector.

We construct such a submatrix iteratively. First, note that any $(n,m,d,k,l,1)$-list UF matrix is also $(n,m,d,k-l,l,1)$-list UF, which  follows from the definition, as any $k-l$ column set can be extended to a $k$ column set by appending $l$ columns.

Let $S \subseteq V$ be any subset of size $l$, for instance
\(
S := \{j_1,...,j_l\}\), and set \( T := V \setminus S.
\)

By the $(n,m,d,k-l,l,1)$-list UF property applied to $(S,T)$, there exist an index $c_1 \in [l]$ and a row index $\lambda_1 \in [m]$ such that
\[
M_{j_{c_1}}[\lambda_1] = 1, \text{ and } M_{j'}[\lambda_1] = 0 \ \text{for all } j'\in \{j_1,...,j_k\} \setminus \{j_{c_1}\}.
\]
Thus, the $\lambda_1$-th row of $M'$ is equal to a standard basis vector, say $e_{i_1}$.

Next, define new subsets
\(
S_2 := \bigl(S \setminus \{j_{c_1}\}\bigr) \cup \{j_{l+1}\},  T_2 := V \setminus S_2.
\)
Applying the same property to $(S_2, T_2)$ yields another row $\lambda_2$ of $M'$ equal to $e_{i_2}$, where $i_2 \neq i_1$ and $\lambda_2\neq \lambda_1$ as equality in either case would result in a collision.

Repeating this procedure, at each step replacing one index in $S$ with a new index from $V \setminus S$, we obtain $k - l + 1$ distinct rows of $M'$, each equal to a distinct standard basis vector.

Therefore, $M'$ contains a submatrix with at least $k - l + 1$ rows forming distinct standard basis vectors, as required.
\end{proof}

\end{document}
